\section{伪标签：从 LLM 打标到规则法替代}
\label{sec:pseudo}

\subsection{Stage 1 的原始设计}

CogECI 的第一阶段要为每个因果正例事件对标注“支撑该因果关系的句子集合”。
按 \texttt{llm/run\_llm\_processing.py} 中的提示词，输入是给每句编号后的整篇文档
与两个事件的触发词及其释义，输出是句子编号列表：

\begin{lstlisting}
The input document is labeled with an ordinal number before each sentence,
e.g. (0,1...) , knowing that event: {e1} and event: {e2} are causally related,
{e1} means [...]. {e2} means [...] identify which sentences from the document
articulate the causal relationship between events {e1} and {e2}. ...
Answer Present the sentence numbers as a list, e.g., "[1,2,...]". don't output other content.
Document :{list(numbered_sentences)}
\end{lstlisting}

注意提示词中\textbf{已告知该事件对确有因果}——Stage 1 不是让 LLM 判断有无因果，
而是让它\emph{找支撑句}。因此只需对 \texttt{target} $\ne 0$ 的正例调用，
数量等于正例数：train 4\,933 次、dev 545 次，合计 5\,478 次；
由于句内流与跨句流使用两个不同变体，实际约 10\,956 次调用。

按真实提示词逐条构造并统计（约 4 字符 = 1 token），单份变体的开销为
输入 4.33\,M token、输出 82\,k token。以 gpt-4o-mini 量级计价约十元人民币，
\textbf{成本不构成障碍}。真正的障碍是调用次数本身：约 1.1 万次串行调用，
按每次 1--3 秒计需 3--9 小时。

\subsection{实际遇到的障碍}

本项目以阿里云百炼 MaaS 的 Qwen3.8-27B 作为 Stage 1 的打标模型，
通过 OpenAI 兼容接口调用。表~\ref{tab:llm} 为尝试过的模型与结果。

\begin{table}[H]
\centering
\caption{LLM 打标模型的实际可用性}
\label{tab:llm}
\small
\begin{tabularx}{\textwidth}{@{}p{4.6cm} l X@{}}
\toprule
模型 & 结果 & 原因\\
\midrule
qwen3.8-27b & dev 全量 545 条完成，train 至 1\,752 次额度耗尽 & 质量最好（strict 平均选 2.19 句）\\
qwen3.8-flash & 同批额度耗尽 & 质量接近 27b\\
deepseek-v4-flash-0731 & 不可用 & 强制思考，\texttt{/no\_think} 无效，思考占满 \texttt{max\_tokens} 导致正文为空\\
qwen3.8-2.4t-a95b & 不可用 & 网关强制 \texttt{enable\_}\allowbreak\texttt{thinking}\allowbreak\texttt{=True}，单次 $>$120\,s 超时\\
kimi-k3\newline qwen3.8-omni-flash & 有额度 & 备用，已在部分 train 文档上使用\\
\bottomrule
\end{tabularx}
\end{table}

其中 qwen3.8-27b 开启思考模式时单次调用需 120\,s 以上、全部超时；
关闭思考（\texttt{enable\_thinking=False}）后降至约 2.6\,s，此时才具备可行性。
deepseek 系列的失败模式值得记录：思考内容会先占满 \texttt{max\_tokens}，
返回的 \texttt{finish\_reason} 为 \texttt{length} 而正文为空，
表现为“83\% 解析失败”；把 \texttt{max\_tokens} 提到 2048 仍不够。

打标进行到 train 的第 1\,752 条时免费额度耗尽，脚本按设计立即中止并保留已写入的 CSV。
此时距离完成 train 全量还差 3\,181 条。训练侧不能无限等待，因此需要离线替代方案。

\subsection{规则法伪标签}

\paragraph{设计原则}伪标签的用途是约束 CGL——让事件对表示靠近“支撑句”、远离其他句。
因此它不需要复现 LLM 的完整推理，只需要给出\emph{高置信度的最小支撑集}。
注意到一个关键事实：因果关系的两个事件各自所在的句子\emph{必然}是支撑句之一，
这是不依赖任何模型的硬约束。据此设计两条规则（记事件 $a,b$ 的所在句编号为
$s_1=\sigma(a)$、$s_2=\sigma(b)$）：

\begin{itemize}
  \item \textbf{strict}（供 \texttt{sen.csv}，即句内流使用）：取 $\{s_1,s_2\}$。
        句内对天然退化为单句。
  \item \textbf{loose}（供 \texttt{sen\_more.csv}，即跨句流使用）：取完整区间
        $[\min(s_1,s_2),\ \max(s_1,s_2)]$；若区间长度超过 8 句，
        退化为 $\{s_1,s_2\}\cup\{s_1\pm1\}\cup\{s_2\pm1\}$ 的局部窗口。
\end{itemize}

\paragraph{为什么要对区间封顶}跨句因果的完整区间在长文档上会失控：
train 的 3\,358 个句间正例中，区间长超过 10 句的占 11.5\%，覆盖超过半篇的占 19.9\%，
最长跨 55 句（等于整篇）。这类标签会让 CGL 的 BCE 目标退化成“全部为 1”，
反而把句向量表示拉平，失去区分度。封顶后，在 dev 上与 LLM loose 答案的
句级 $F_1$ 由 0.723 提升到 \textbf{0.746}，标签密度也更接近 LLM 的实际输出
（规则 2.77 句/对 vs.\ LLM 3.13 句/对）。

\subsection{一致性对照实验}

规则法是否可信，不能靠主观判断，必须与 LLM 标注做定量比对。
本项目的有利条件是：dev 的 545 条 LLM 标注已全部完成，可作为参照集。
在 dev 全量正例上按句级 micro $P/R/F_1$ 统计，结果见表~\ref{tab:agree}。

\begin{table}[H]
\centering
\caption{规则法伪标签与 Qwen3.8-27B 标注的一致性（dev，545 个正例）}
\label{tab:agree}
\footnotesize
\setlength{\tabcolsep}{3.2pt}
\begin{tabular}{lcccccc}
\toprule
对比 & 句级 $P$ & 句级 $R$ & 句级 $F_1$ & pair 级 $F_1$
  & 规则集 $\subseteq$ LLM 集 & LLM 含两事件句\\
\midrule
规则 strict vs.\ LLM strict & \textbf{0.985} & 0.753 & 0.853 & 0.891 & 97.4\% & \textbf{97.4\%}\\
规则 strict vs.\ LLM loose  & 0.998 & 0.531 & 0.693 & 0.756 & 99.6\% & 99.6\%\\
\bottomrule
\end{tabular}
\end{table}

三个数字值得展开：

\begin{enumerate}
  \item \textbf{精确率 0.985}：规则选出的句子有 98.5\% 也被 LLM 选中。
        这意味着规则法几乎不引入噪声——对 CGL 这类“用伪标签约束表示”的弱监督而言，
        精度比召回更重要，因为错误的负样本会把真实支撑句推远。
  \item \textbf{“LLM 含两个事件句” 达 97.4\%}：这是不依赖任何规则设计的硬指标，
        只问“LLM 的答案里是否包含两个事件所在句”。97.4\% 说明
        \textbf{规则法捕捉的正是 LLM 选句的核心}，而非巧合。
  \item \textbf{召回 0.753 偏低}：LLM 平均会多选约 0.5 句，典型情况是把标题句也算作支撑，
        或多带一句因果承接句。这部分属于“可选上下文”，不是“必需上下文”。
\end{enumerate}

进一步按事件对类型拆分，用于在多个候选规则中选定最终方案，见表~\ref{tab:reltype}。

\begin{table}[H]
\centering
\caption{按事件对类型拆分的规则对比（dev）}
\label{tab:reltype}
\small
\begin{tabular}{llcrccc}
\toprule
分组 & 参照 & 规则 & $n$ & $P$ & $R$ & $F_1$\\
\midrule
intra & LLM strict & A $=\{s_1,s_2\}$ & 176 & \textbf{0.994} & 0.764 & 0.864\\
inter & LLM strict & A $=\{s_1,s_2\}$ & 369 & 0.982 & 0.751 & 0.851\\
\midrule
inter & LLM loose & B $=$ 完整区间 & 329 & 0.690 & 0.759 & 0.723\\
inter & LLM loose & C $=$ 区间 $\pm1$ & 329 & 0.539 & 0.840 & 0.657\\
inter & LLM loose & \textbf{D $=$ 区间封顶 8 $\pm1$} & 329 & \textbf{0.748} & 0.745 & \textbf{0.746}\\
\bottomrule
\end{tabular}
\end{table}

最终选定 \textbf{intra 用 A、inter 用 D}，即前文的 strict / loose 两个变体。
规则 C 的召回最高（0.840）但精度只有 0.539，说明盲目扩窗会引入大量非支撑句，故弃用。

\subsection{产物与校验}

四份伪标签文件（train/dev $\times$ strict/loose）经终检全部通过：

\begin{table}[H]
\centering
\caption{伪标签文件终检结果}
\label{tab:validate}
\small
\begin{tabular}{lccc}
\toprule
文件 & 覆盖文档 & \texttt{answer} 总数 & 正例缺失\\
\midrule
\texttt{llama3help\_train\_sen.csv} & 233/233 & 61\,065 & 0/4\,933\\
\texttt{llama3help\_train\_sen\_more.csv} & 233/233 & 61\,065 & 0/4\,933\\
\texttt{llama3help\_dev\_sen.csv} & 25/25 & 9\,514 & 0/545\\
\texttt{llama3help\_dev\_sen\_more.csv} & 25/25 & 9\,514 & 0/545\\
\bottomrule
\end{tabular}
\end{table}

此外，全部 LLM 调用的完整提示词与原始响应都写入了 JSONL 留痕文件
（\texttt{raw\_*.jsonl}），满足课程“保留提示与原始响应”的要求。

\paragraph{合规性说明}本方案不触碰测试标签：伪标签按文档生成；5 折交叉验证中，
测试折的文档不参与训练，其伪标签不产生梯度；参与梯度更新的只有训练折文档，
而这些标签来自文档\emph{自身}的 gold 因果标注——这正是原论文
“label-informed pseudo labelling”的设计。差异在于：本次伪标签由规则而非 LLM 生成，
这一点已在第~\ref{sec:discussion} 节列为局限。
